\section{Relèvement nonadique de l'unité de Pell}

Pour \(k\ge1\), soit l'anneau quadratique non ramifié
\begin{equation}
\mathcal R_k=(\Z/3^k\Z)[r]/(r^2-2),
\label{eq:pell-ring}
\end{equation}
et soit

\[
u=3+2r.
\]

La conjugaison \(r\mapsto-r\) donne
\[
N_{\mathcal R_k/(\Z/3^k\Z)}(u)
=(3+2r)(3-2r)=9-8=1,
\]
de sorte que \(u\) est une unité de norme un. En outre,
\[
u^2=17+12r,\qquad
u^4-1=576+408r=3(192+136r).
\]
Le second facteur est une unité modulo trois parce qu'il se réduit à \(r\ne0\). Par
conséquent, la valuation \(3\)-adique satisfait
\(v_3(u^4-1)=1\). Le lemme de relèvement pour une unité congrue à un
modulo trois produit, pour tout \(j\ge0\),
\[
v_3\!\left(u^{4\cdot3^j}-1\right)=j+1.
\]
Comme \(u\bmod3=2r\) et \((2r)^2=-1\), son ordre modulo trois est quatre.
En conséquence,

\begin{equation}
\operatorname{ord}(u)=4\cdot3^{k-1},
\label{eq:pell-order}
\end{equation}

et la clôture profinie des groupes cycliques compatibles est
\begin{equation}
\varprojlim_k\langle u\bmod3^k\rangle
\simeq C_4\times\mathbb Z_3.
\label{eq:pell-profinite}
\end{equation}
Sa réalisation archimédienne est

\[
3+2\sqrt2=(1+\sqrt2)^2.
\]

Cette construction démontre la conservation de la profondeur nonadique
par une unité hyperbolique. Les unités \(3+2\sqrt2\),
\(2+\sqrt3\) et \(30+\sqrt{899}\) forment une famille de Pell de
provenances distinctes ; toute identification entre elles exige des
applications explicites entre les extensions correspondantes.

